% (User input window)

\paragraph{两温度失败指数的变分结构：幂律耦合、主序链与临界二次律}
本段把定理 \ref{thm:pom-microcanonical-hypergeometric-ldp-kl-sym} 的 KL--对称速率函数
从 \((q\mapsto J_\beta(q))\) 的单变量表述提升为一个两温度变分问题，并在熵阈值的超临界侧给出唯一极小化器的结构。
其结果可作为“剩余组成熵偏大”事件的指数级失败证书，并与定理 \ref{thm:pom-microcanonical-query-plus-bits-phase} 的临界窗叙述兼容。

\begin{definition}[两温度失败指数：熵阈值的最小 KL--对称代价]\label{def:pom-microcanonical-two-temperature-failure-exponent}
设 \(X\) 为有限集（\(|X|=n\ge 2\)），取全支撑分布 \(w\in\Delta(X)\)，并固定 \(\beta\in(0,1)\)。
对任意 \(q,r\in\Delta(X)\) 满足质量守恒
\[
w=\beta q+(1-\beta)r,
\]
定义两温度散度泛函
\[
J_\beta(q,r):=\beta D(q\|w)+(1-\beta)D(r\|w).
\]
对熵阈值 \(h>H(w)\)，定义超临界失败指数
\[
E_{\mathrm{fail}}(\beta,h)
:=
\inf\Bigl\{J_\beta(q,r):\ q,r\in\Delta(X),\ w=\beta q+(1-\beta)r,\ H(r)\ge h\Bigr\}.
\]
若可行域为空，则约定 \(E_{\mathrm{fail}}(\beta,h)=+\infty\)。
\end{definition}

\begin{theorem}[两温度 KKT 结构：唯一极小化器的幂律耦合]\label{thm:pom-microcanonical-two-temperature-kkt-power-law}
\leanverified{paper\_pom\_microcanonical\_two\_temperature\_kkt\_power\_law}
沿用定义 \ref{def:pom-microcanonical-two-temperature-failure-exponent}，并假设给定 \(h\) 使可行域非空且 \(h>\!H(w)\)。
则极小化问题存在唯一极小点 \((q^\star,r^\star)\)，并且熵约束必饱和：
\[
H(r^\star)=h.
\]
此外，存在唯一乘子 \(\eta^\star>0\)（因而存在唯一指数 \(\gamma>1\)）使得对所有 \(x\in X\),
\[
q^\star(x)=\frac{(r^\star(x))^{\gamma}}{\sum_{y\in X}(r^\star(y))^{\gamma}},
\qquad
\gamma=1+\frac{\eta^\star}{1-\beta},
\]
等价地
\(
\log q^\star=\gamma\log r^\star+\mathrm{const}
\)
（常数在 \(x\) 上不依赖）。
\end{theorem}
\begin{proof}
可行域由线性等式与凹函数 \(H\) 的上水平集给出，故为紧凸集；而 \(J_\beta\) 作为两项 KL 散度的正权和，在 \(\Delta(X)^2\) 的内点上严格凸，从而极小点唯一。
\par
由于 \(w\) 全支撑且 \(\beta\in(0,1)\)，若 \((q,r)\) 可行，则 \(q,r\) 自动全支撑；因此可在内点上写 KKT 系统。
若熵约束不饱和，则对应乘子 \(\eta^\star=0\)，从而极小点必须同时极小化仅含质量守恒的松弛问题；而该松弛问题的唯一极小点为 \(q=r=w\)，与 \(h>H(w)\) 矛盾。
故 \(\eta^\star>0\) 且 \(H(r^\star)=h\)。
\par
对约束写拉格朗日函数并对每个坐标取驻点条件，消去质量守恒的乘子后可得
\(
\log r^\star=\alpha\log q^\star+\mathrm{const}
\),
其中
\(
\alpha=\frac{1-\beta}{1-\beta+\eta^\star}\in(0,1)
\)。
从而 \(q^\star\propto (r^\star)^{1/\alpha}\)，并令 \(\gamma:=1/\alpha=1+\eta^\star/(1-\beta)\) 得到所示幂律耦合。
唯一性迫使 \(\eta^\star\) 与 \(\gamma\) 唯一。证毕。
\end{proof}

\begin{proposition}[幂律轨道的存在唯一性与熵的严格单调性]\label{prop:pom-microcanonical-two-temperature-escort-orbit-entropy-monotone}
沿用定义 \ref{def:pom-microcanonical-two-temperature-failure-exponent}。
对每个 \(\gamma>1\)，定义幂 escort 变换
\[
\mathrm{Escort}_\gamma(r)(x):=\frac{r(x)^\gamma}{\sum_{y\in X}r(y)^\gamma}\qquad(r\in\Delta(X)).
\]
若存在 \((q,r)\in\Delta(X)^2\) 满足 \(w=\beta q+(1-\beta)r\) 且 \(q=\mathrm{Escort}_\gamma(r)\)，则该解唯一；记之为 \((q_\gamma,r_\gamma)\)。
并且在其定义域上，映射
\[
\gamma\longmapsto H(r_\gamma)
\]
连续且严格递增，并满足
\[
\lim_{\gamma\downarrow 1}r_\gamma=w.
\]
令
\[
H_{\max}(\beta,w):=\sup\Bigl\{H(r):\ r\in\Delta(X),\ \frac{w-(1-\beta)r}{\beta}\in\Delta(X)\Bigr\}\in(H(w),\log n],
\]
则 \(\lim_{\gamma\to\infty}H(r_\gamma)=H_{\max}(\beta,w)\)。
特别地，若均匀分布 \(U\) 可行（等价于 \(w(x)\ge (1-\beta)/n\) 对所有 \(x\) 成立），则 \(H_{\max}(\beta,w)=\log n\) 且 \(r_\gamma\to U\)。
因此对任意 \(h\in(H(w),H_{\max}(\beta,w))\) 存在唯一 \(\gamma(h)>1\) 使 \(H(r_{\gamma(h)})=h\)，并且该 \((q_{\gamma(h)},r_{\gamma(h)})\) 与定理 \ref{thm:pom-microcanonical-two-temperature-kkt-power-law} 的极小化器一致。
\end{proposition}
\begin{proof}
对固定 \(\eta>0\) 考虑惩罚型问题
\[
\mathcal{V}(\eta):=\inf\Bigl\{J_\beta(q,r)-\eta H(r):\ q,r\in\Delta(X),\ w=\beta q+(1-\beta)r\Bigr\}.
\]
由于 \(-H\) 为严格凸函数且 \(J_\beta\) 严格凸，故该问题存在唯一极小点 \((q_\eta,r_\eta)\)；
其 KKT 系统给出幂律关系 \(q_\eta=\mathrm{Escort}_\gamma(r_\eta)\) 且 \(\gamma=1+\eta/(1-\beta)>1\)。
因此对每个 \(\gamma>1\) 至多对应一个可行解，记为 \((q_\gamma,r_\gamma)\)。
\par
函数 \(\eta\mapsto \mathcal{V}(\eta)\) 为 \(\eta\) 的下确界包络，因而为凹函数；并且在唯一极小点下可微且满足
\(
\mathcal{V}'(\eta)=-H(r_\eta)
\)。
凹性推出 \(-\mathcal{V}'(\eta)=H(r_\eta)\) 随 \(\eta\)（从而随 \(\gamma\)）严格递增，连续性由隐函数/稳定性与唯一极小性给出。
\par
当 \(\eta\downarrow 0\) 时，惩罚项消失且松弛问题的唯一极小点为 \(q=r=w\)，故 \(r_\gamma\to w\)（\(\gamma\downarrow 1\)）。
当 \(\eta\to\infty\) 时，\(-\eta H(r)\) 主导并迫使 \(r_\eta\) 收敛到“在可行多面体内最大化 \(H\)”的唯一解，从而 \(H(r_\gamma)\uparrow H_{\max}(\beta,w)\)；
若 \(U\) 可行，则最大熵点即为 \(U\) 且极限为 \(\log n\)。
最后，对 \(h\in(H(w),H_{\max})\) 由严格单调性得到唯一 \(\gamma(h)\)，并与定理 \ref{thm:pom-microcanonical-two-temperature-kkt-power-law} 的唯一极小化器匹配。证毕。
\end{proof}

\begin{corollary}[共同排序与主序链：\texorpdfstring{$r^\star\prec w\prec q^\star$}{r* ≺ w ≺ q*}]\label{cor:pom-microcanonical-two-temperature-majorization-chain}
在定理 \ref{thm:pom-microcanonical-two-temperature-kkt-power-law} 的设定下，存在同一排列使得 \(r^\star,w,q^\star\) 的坐标按非增序排列后满足逐段和链
\[
\sum_{i=1}^k r^\star_{(i)}\ \le\ \sum_{i=1}^k w_{(i)}\ \le\ \sum_{i=1}^k q^\star_{(i)}\qquad(k=1,\dots,n-1),
\]
且至少对某个 \(k\) 严格不等；从而 \(r^\star\prec w\prec q^\star\)。
\end{corollary}
\begin{proof}
由幂律耦合 \(q^\star=\mathrm{Escort}_\gamma(r^\star)\)（\(\gamma>1\)），两者具有相同坐标排序，且幂变换在主序意义上将分布严格冷化：\(r^\star\prec q^\star\)（Hardy--Littlewood--Pólya）。
又因 \(w=\beta q^\star+(1-\beta)r^\star\) 为同序向量的凸组合，逐段和对凸组合线性，故得到链式不等式并推出 \(r^\star\prec w\prec q^\star\)。
严格性由 \(h>H(w)\) 排除 \(r^\star=w\)（从而亦排除 \(q^\star=w\)）得到。证毕。
\end{proof}

\begin{proposition}[Jensen--Shannon 结构恒等式：\texorpdfstring{$J_\beta$}{Jbeta} 等于混合熵亏损]\label{prop:pom-microcanonical-two-temperature-js-identity}
\leanverified{paper\_pom\_microcanonical\_two\_temperature\_js\_identity}
在约束 \(w=\beta q+(1-\beta)r\) 下，对任意可行 \((q,r)\) 有恒等式
\[
J_\beta(q,r)=H(w)-\beta H(q)-(1-\beta)H(r).
\]
因此当可行域非空且 \(h>H(w)\) 时，
\[
E_{\mathrm{fail}}(\beta,h)
=H(w)-(1-\beta)h-\beta\sup\Bigl\{H(q):\exists r,\ w=\beta q+(1-\beta)r,\ H(r)=h\Bigr\}.
\]
\end{proposition}
\begin{proof}
由 \(D(p\|w)=-H(p)-\sum_x p(x)\log w(x)\)，并用质量守恒
\(
\beta q+(1-\beta)r=w
\)
消去 \(\sum_x(\beta q+(1-\beta)r)\log w\)，即得第一式。
第二式由定理 \ref{thm:pom-microcanonical-two-temperature-kkt-power-law} 的 \(H(r^\star)=h\) 代入并对 \(q\) 作等价极大化得到。证毕。
\end{proof}

\begin{theorem}[强对偶、唯一乘子与敏感性公式]\label{thm:pom-microcanonical-two-temperature-strong-duality-sensitivity}
\leanverified{paper\_pom\_microcanonical\_two\_temperature\_strong\_duality\_sensitivity}
沿用定义 \ref{def:pom-microcanonical-two-temperature-failure-exponent}，假设给定 \(h\in(H(w),H_{\max}(\beta,w))\) 使可行域非空。
定义对偶值函数（\(\eta\ge 0\)）
\[
\mathcal{V}(\eta):=\inf\Bigl\{J_\beta(q,r)-\eta H(r):\ q,r\in\Delta(X),\ w=\beta q+(1-\beta)r\Bigr\}.
\]
则成立强对偶
\[
E_{\mathrm{fail}}(\beta,h)=\sup_{\eta\ge 0}\bigl(\mathcal{V}(\eta)+\eta h\bigr),
\]
且上确界在唯一 \(\eta^\star>0\) 处达到。
并且函数 \(h\mapsto E_{\mathrm{fail}}(\beta,h)\) 在 \((H(w),H_{\max}(\beta,w))\) 上可微，且
\[
\frac{\partial}{\partial h}E_{\mathrm{fail}}(\beta,h)=\eta^\star,
\qquad
\gamma(h)=1+\frac{\eta^\star}{1-\beta},
\]
其中 \(\gamma(h)\) 为定理 \ref{thm:pom-microcanonical-two-temperature-kkt-power-law} 的唯一幂指数。
\end{theorem}
\begin{proof}
约束 \(H(r)\ge h\) 等价于凸不等式 \(h-H(r)\le 0\)，并且在 \(h< H_{\max}(\beta,w)\) 时满足 Slater 条件（存在严格可行点）。
因此强对偶成立，且由于原问题极小点唯一，KKT 乘子亦唯一且满足 \(\eta^\star>0\)（同定理 \ref{thm:pom-microcanonical-two-temperature-kkt-power-law}）。
可微性与导数公式由凸对偶的一般敏感性定理（包络定理）推出；
最后一式由幂律关系 \(\gamma=1+\eta/(1-\beta)\) 代入得到。证毕。
\end{proof}

\begin{theorem}[临界熵窗的二次律：\texorpdfstring{$h\downarrow H(w)$}{h↓H(w)} 时的精确起跳系数]\label{thm:pom-microcanonical-two-temperature-critical-quadratic-law}
设 \(w\) 全支撑且非均匀，并记
\[
V(w):=\Var_{X\sim w}\bigl(\log w(X)\bigr)>0
\]
（同推论 \ref{cor:pom-microcanonical-information-clt}）。
令 \(\delta:=h-H(w)\downarrow 0^+\) 且取 \(h\in(H(w),H_{\max}(\beta,w))\)。
则有渐近展开
\[
E_{\mathrm{fail}}(\beta,h)
=\frac{1-\beta}{2\beta}\cdot \frac{\delta^2}{V(w)}+o(\delta^2).
\]
若改用比特密度
\(
b:=\frac{(1-\beta)h}{\log 2}
\)
与临界点
\(
b_{\mathrm{crit}}:=\frac{(1-\beta)H(w)}{\log 2}
\),
则等价地
\[
E_{\mathrm{fail}}(\beta,b)
=\frac{(b-b_{\mathrm{crit}})^2(\log 2)^2}{2\beta(1-\beta)V(w)}+o\!\left((b-b_{\mathrm{crit}})^2\right).
\]
\end{theorem}
\begin{proof}
当 \(\delta\downarrow 0\) 时，最优点 \((q^\star,r^\star)\) 逼近 \((w,w)\)。
写 \(r=w+u\)（\(\sum_x u(x)=0\)），则由质量守恒 \(q=w-\frac{1-\beta}{\beta}u\)。
对 KL 在 \(w\) 处作二阶展开可得
\[
J_\beta(q,r)
=\frac{1-\beta}{2\beta}\sum_{x\in X}\frac{u(x)^2}{w(x)}+o(\|u\|_2^2).
\]
另一方面，熵的一阶变分（利用 \(\sum_x u(x)=0\)）给出
\[
H(w+u)-H(w)=-\sum_{x\in X}u(x)\log w(x)+o(\|u\|_2),
\]
故主阶约束为 \(\sum_x u(x)\log w(x)=-\delta\)。
因此主阶问题化为：在二次型 \(\sum u^2/w\) 下最小化线性约束 \(\sum u\log w=-\delta\)。
拉格朗日乘子法给出唯一主方向
\[
u^\star(x)=-\frac{\delta}{V(w)}\,w(x)\Bigl(\log w(x)-\EE_w[\log w]\Bigr),
\]
并满足
\(
\sum_x\frac{(u^\star(x))^2}{w(x)}=\delta^2/V(w)
\)。
代回即得
\(
E_{\mathrm{fail}}(\beta,h)=\frac{1-\beta}{2\beta}\cdot \frac{\delta^2}{V(w)}+o(\delta^2)
\)。
比特密度形式由 \(\delta=\frac{\log 2}{1-\beta}(b-b_{\mathrm{crit}})\) 代换得到。证毕。
\end{proof}
